%% kapa: exam=1 round=25 year=2014 period=2 form=A subject=appraisal_law subject_ko=부동산관계법규 seq=39 number=39
%% source: Q-Net 감정평가사 기출문제 (cst00309) artlSeq=5101207 file=extracted/1008/5101207/2014년 제25회 감정평가사 1차시험 기출문제/2교시(A형).hwp
%% conversion: kordoc 4.24.51 -> tools/gen1.py (auto); review=proofread-sonnet; flags=none
\begin{question}{39}
부동산등기법령상 가등기에 관한 설명으로 옳지 않은 것은?
\begin{choices}
  \item 가등기권리자는 가등기의무자의 승낙이 있을 때에는 단독으로 가등기를 신청할 수 있다.
  \item 가등기에 의한 본등기를 한 경우 본등기의 순위는 가등기의 순위에 따른다.
  \item 가등기명의인은 단독으로 가등기의 말소를 신청할 수 있다.
  \item 가등기에 의한 본등기가 이루어졌을 때 가등기 이후에 된 등기로서 가등기에 의하여 보전되는 권리를 침해하는 등기에 대해서는 등기상 이해관계인의 신청에 의하여 등기관이 이를 말소한다.
  \item 소유권 이전을 위한 정지조건부 청구권을 보전하기 위한 가등기도 가능하다.
\end{choices}
\end{question}
\answer{4}
